\subsection{Module 0: Developmental State, Excitable Element and Refractory Clock}\label{sec:m0}

\paragraph{Role.}
Before collective waves start, isolated starved cells emit cAMP pulses spontaneously, on average \mk{1}{about once per 12\,min}~\cite{ford2023} and \mk{2}{at most about once per 4\,min}~\cite{ford2023}, with inter-pulse intervals that \mk{3}{follow a gamma distribution (shape 6.6, rate {$1.25\,\mathrm{min^{-1}}$})}~\cite{ford2023}, i.e.\ \mk{5}{a mean of 5.3\,min}~\cite{ford2023}.
Such \mk{9}{noise-driven pulsing of single cells below the threshold for collective behaviour is a key step in the onset of oscillations}~\cite{gregor2010,sgro2015}.
The molecular circuit that generates periodic cAMP in populations couples \mk{5}{cAMP-induced activation of ACA and ERK2}~\cite{laub1998,maeda2004}, \mk{6}{inhibition of ERK2 by PKA}~\cite{laub1998,maeda2004} and \mk{7}{phosphorylation of the phosphodiesterase RegA by ERK2}~\cite{laub1998,maeda2004}, and produces a period of \mk{8}{about 7\,min}~\cite{laub1998,maeda2004}.
Module~0 replaces the unresolved part of this circuit by a compact excitable element and adds the two variables that the collective behaviour needs: a developmental state and a refractory clock.

\paragraph{Why the driver is X and not cAMP.}
In the literature circuit cAMP is the input that activates ACA and ERK2.
In the reactions of this module (Table~\ref{tab:m0rxn}) ERK2 and ACA are instead activated by the abstract activator X (0.9e, 0.9f); this is a deliberate design choice and not a measured dependence.
Wiring cAMP into these steps does not work in this network, for two reasons found while building the model.
First, every route from the receptor to the cyclase passes through receptor occupancy, so an isolated cell without external cAMP has no input to ACA at all and can never fire, although isolated starved cells do pulse spontaneously~\cite{ford2023}.
Second, a feedback of intracellular cAMP onto ACA cannot produce a pulse: all paths from cAMP$_i$ to ACA run through PKA, whose edges are inhibitory, and an added cooperative autocatalysis on cAMP$_i$ either settles on a plateau or latches in the high state; neither has a reset.
A pulse needs a fast autocatalysis together with a slow recovery, and the receptor-independent arm of the network has neither.
X supplies both (Schl\"ogl autocatalysis with the recovery variable Y, a FitzHugh--Nagumo pair), is seeded by the Hunger-gated ignition 0.9a so that it needs no receptor input, and is therefore the node at which an isolated cell decides to fire.
X thus stands for the unresolved part of the circuit between cAMP and ACA/ERK2 and is not a molecule that has been identified.
Consequently the model contains no direct activation of ACA or ERK2 by cAMP in the table of this module; the cAMP-dependence of the circuit enters only through the relay that fires X (Module~8).

\begin{table*}[!t]
\centering
\caption{Module~0 reactions: Hunger, the excitable element (X, Y), the refractory clock (Rf) and the ERK2$\to$ACA output stage. For the species of the excitable element $N_s$ is the copy number of species $s$ \emph{per cell}; $c(\mathrm H)$ and $b(\mathrm{Rf})$ are the Hill gates of Sec.~\ref{sec:hill}.}
\label{tab:m0rxn}
\footnotesize
\renewcommand{\arraystretch}{1.15}
\begin{tabularx}{\textwidth}{@{}l>{\raggedright\arraybackslash}p{0.37\textwidth}>{\raggedright\arraybackslash}X@{}}
\toprule
ID & Reaction & Propensity (rate law)\\
\midrule
0.8a & $\varnothing \to \mathrm{H}$ & $k_H\,H_{\max}\ \text{(zeroth order)}$ \\
0.8b & $\mathrm{H} \to \varnothing$ & $k_H\, N_\mathrm{H}$ \\
0.9a & $\varnothing \to \mathrm{X}$ & $k_0\, c(\mathrm{H})\, b(\mathrm{Rf})\quad[\text{molec\,s}^{-1}\text{cell}^{-1}]$ \\
0.9b & $2\mathrm{X} \to 3\mathrm{X}$ & $\tfrac{k_1}{2}N_\mathrm{X}(N_\mathrm{X}-1)$ \\
0.9c & $3\mathrm{X} \to 2\mathrm{X}$ & $\tfrac{k_6}{6}N_\mathrm{X}(N_\mathrm{X}-1)(N_\mathrm{X}-2)$ \\
0.9d & $\mathrm{X} \to \varnothing$ & $k_2 N_\mathrm{X}$ \\
0.9y1 & $2\mathrm{X} \to 2\mathrm{X}+\mathrm{Y}$ & $\tfrac{k_3}{2}N_\mathrm{X}(N_\mathrm{X}-1)$ \\
0.9y2 & $\mathrm{X}+\mathrm{Y} \to \mathrm{Y}$ & $k_5 N_\mathrm{X}N_\mathrm{Y}$ \\
0.9y3 & $\mathrm{Y} \to \varnothing$ & $k_4 N_\mathrm{Y}$ \\
0.9r1 & $2\mathrm{X} \to 2\mathrm{X}+\mathrm{Rf}$ & $\tfrac{k_{Rp}}{2}N_\mathrm{X}(N_\mathrm{X}-1)$ \\
0.9r2 & $\mathrm{Rf} \to \varnothing$ & $k_{Rd} N_\mathrm{Rf}$ \\
0.9e & $2\mathrm{X}+\mathrm{ERK2} \to 2\mathrm{X}+\mathrm{ERK2}^*$ & $k_{XE}\,[\mathrm{X}]^2[\mathrm{ERK2}]\ \text{(falling factorial in }N_\mathrm{X})$ \\
0.9e$'$ & $\mathrm{ERK2}^* \to \mathrm{ERK2}$ & $k_{Ed}[\mathrm{ERK2}^*]$ \\
0.9f & $\mathrm{X}+\mathrm{ACA} \to \mathrm{X}+\mathrm{ACA}^*$ & $k_{XA}[\mathrm{X}][\mathrm{ACA}]\,\dfrac{[\mathrm{ERK2}^*]^4}{K_E^4+[\mathrm{ERK2}^*]^4}$ \\
\bottomrule
\end{tabularx}
\end{table*}

\begin{table*}[!t]
\centering
\caption{Module~0 constants. Basis: L literature value, D derived from a literature anchor, C calibrated on the model against a stated target, E estimate or design choice (Sec.~\ref{sec:conv}).}
\label{tab:m0const}
\footnotesize
\renewcommand{\arraystretch}{1.15}
\begin{tabularx}{\textwidth}{@{}l r l c >{\raggedright\arraybackslash}X@{}}
\toprule
Symbol & Value & Unit & Basis & Meaning, justification, source\\
\midrule
$H_{\max}$ & \num{100} & $\mathrm{molec/voxel}$ & E & Steady-state level of the abstract \emph{Hunger} variable; held at full developmental competence. \\
$k_H$ & \num{2.8e-4} & $\mathrm{s^{-1}}$ & E & Relaxation of Hunger ($\tau\approx 1$\,h); the production rate is $H_{\max}k_H$. A slow developmental clock, not a measured rate. \\
$H_{1/2},\ n_c$ & \num{60},\ 4 & $\mathrm{molec/voxel}$ & E & Half-competence level and steepness of the competence gate; $c(H_{\max})=0.885$. \\
$k_0$ & \num{3.5} & $\mathrm{molec\,s^{-1}\,cell^{-1}}$ & C & Ignition (seed) rate at full competence. Set on the stochastic engine so that the mean inter-pulse interval of an isolated cell is 726\,s; target \mkt{1}{{$\approx\!12$}\,min}~\cite{ford2023}. \\
$k_1$ & \num{0.08} & $\mathrm{molec^{-1}s^{-1}}$ & C & Autocatalysis ($2\mathrm X\to3\mathrm X$); creates the excitation threshold ($\approx\!23$ molecules). \\
$k_6$ & \num{5e-4} & $\mathrm{molec^{-2}s^{-1}}$ & C & Cubic self-limitation; bounds the spike at $\approx\!456$ molecules per cell. \\
$k_2$ & \num{1} & $\mathrm{s^{-1}}$ & C & Basal removal of X (relaxation time 1\,s). \\
$k_3$ & \num{6.25e-5} & $\mathrm{molec^{-1}s^{-1}}$ & C & Spike-triggered production of the recovery variable Y. \\
$k_5$ & \num{0.05} & $\mathrm{molec^{-1}s^{-1}}$ & C & Y-mediated spike termination. \\
$k_4$ & \num{0.01} & $\mathrm{s^{-1}}$ & C & Decay of Y ($\tau=100$\,s); sets the recovery of the element. \\
$k_{Rp},\ k_{Rd}$ & \num{5e-3},\ \num{5e-3} & $\mathrm{molec^{-1}s^{-1}},\ \mathrm{s^{-1}}$ & C & Charging and decay ($\tau_R=200$\,s) of the refractory clock. \\
$K_R,\ n_R$ & \num{450},\ 4 & $\mathrm{molec/cell}$ & C & Half-block level and steepness of the refractory gate (Sec.~\ref{sec:hill}). \\
$k_{XE}$ & \num{1823} & $\mu\mathrm{M}^{-2}\mathrm{s}^{-1}$ & C & X$^2$-driven activation of ERK2 (cell units); second order in X to raise the contrast between spike and resting leak. \\
$k_{Ed}$ & \num{0.02} & $\mathrm{s^{-1}}$ & E & ERK2 dephosphorylation ($\tau=50$\,s); ERK2 activation is \mkt{1}{transient}~\cite{maeda1996}. \\
$k_{XA}$ & \num{65} & $\mu\mathrm{M}^{-1}\mathrm{s}^{-1}$ & C & X-driven ACA activation, gated on ERK2$^*$; sets the $\approx\!10$\,s rise of ACA$^*$. \\
$K_E,\ n_E$ & \num{192},\ 4 & $\mathrm{molec/voxel}$ & C & ERK2$^*$ level (12.5\,\% of the pool) at half-open ACA gate; enforces the ordering ERK2$\to$ACA. \\
\midrule
$N_\mathrm{vox}$ & \num{13} &  & E & Voxels per cell (plus-shaped footprint of $52\,\mu\mathrm{m}^2$, i.e.\ about 8\,$\mu$m diameter, in a 2\,$\mu$m thick layer, $104\,\mu\mathrm{m}^3$); real cells have a \mkt{4}{diameter of about 10\,{$\mu$}m}~\cite{ford2023}, so the cell volume is underestimated. $M_c=13\times4817.6=62\,629$ molecules\,$\mu\mathrm{M}^{-1}$ per cell. \\
$[\mathrm{ERK2}]_\mathrm{tot}$ & \num{1540} & $\mathrm{molec/voxel}$ & E & Pool size on the scale of the other kinases ($0.32\,\mu$M). \\
$[\mathrm{ACA}]_\mathrm{tot}$ & \num{865} & $\mathrm{molec/voxel}$ & C & Chosen so that one pulse exports $\approx\!10^7$ cAMP molecules per cell at $k_9+k_{11}=20\,\mathrm{s^{-1}}$ (Module~9). \\
$D_\mathrm{X},D_\mathrm{Y},D_\mathrm{Rf},D_\mathrm{ERK2}$ & \num{10} & $\mu\mathrm{m^2/s}$ & E & Generic cytosolic value; the cell is mixed in $\approx\!2.5$\,s, which matters because the gates are evaluated per voxel. \\
\bottomrule
\end{tabularx}
\end{table*}


\paragraph{Reactions.}
\begin{description}\itemsep0pt
\item[0.8a/b] An abstract \emph{Hunger} variable $H$ represents the developmental state of a starved cell; it is not a molecule. It scales the competence to fire (gate $c(H)$), the cAMP synthesis rate (9.1) and the induction of the extracellular phosphodiesterase system (Module~11). In the simulations $H$ starts at its steady state $H_{\max}$ and fluctuates around it, so it acts as an approximately constant developmental level.
\item[0.9a] Spontaneous ignition: a rare, competence-gated seed that injects X, the activator, at $k_0=3.5$ molecules\,s$^{-1}$ per cell. It is switched off by the refractory clock, which makes the block hard: with the seed shut, X decays through 0.9d and the autocatalysis 0.9b, which needs $N_\mathrm{X}\ge2$, has nothing to amplify.
\item[0.9b--d] The autocatalytic step 0.9b creates an excitation threshold, the cubic loss 0.9c bounds the spike and keeps the stochastic simulation finite, and the linear loss 0.9d sets the resting state. These three reactions are the \mk{4}{chemical Schl\"ogl model}~\cite{schlogl1972} of a \mk{3}{bistable}~\cite{schlogl1972}/excitable system, in a form in which the buffered reservoirs A, B and C of the original scheme (A${}+2$X$\rightleftharpoons3$X, B${}+$X$\rightleftharpoons$C) are absorbed into the rate constants: 0.9b is the forward and 0.9c the reverse direction of A${}+2$X$\to3$X (the $k_1an^2-k_1'n^3$ of the original rate law), and 0.9d is B${}+$X$\to$C. The reverse direction C$\to$B${}+$X is not a separate reaction; its role as the source of X is taken by the gated ignition 0.9a. The propensities are the stochastic (falling-factorial) form of the Schl\"ogl rate law.
\item[0.9y1--3] The recovery variable Y is produced only by spikes (quadratic in X, so resting fluctuations do not raise it), terminates the spike by removing X (0.9y2) and decays slowly (0.9y3). X and Y form the chemical counterpart of the \mk{2}{FitzHugh--Nagumo excitable system with activator and delayed inhibitor}~\cite{fitzhugh1961,nagumo1962}. With the constants of Table~\ref{tab:m0const} and the competence gate at $H_{\max}$ ($c=0.885$), the cell-level nullcline gives a resting level $\approx3.4$, an excitation threshold $\approx24$ and a spike maximum $\approx456$ molecules per cell.
\item[0.9r1/2] A refractory clock Rf is charged by spikes of X (quadratic, as for Y) and decays with $\tau_R=200$\,s. Y cannot serve as the clock because it also terminates the spike, and its time constant would then be tied to the pulse shape. Charging the clock from X makes the block \emph{source independent}: a cell that has fired, spontaneously or after a relayed signal (Module~8), cannot be fired again by either route until the clock has decayed. The model requires a refractory period of several minutes, commensurate with the \mk{8}{$\approx$7\,min oscillation period}~\cite{maeda2004} and the \mk{2}{$\approx$4\,min minimum interval between pulses of single cells}~\cite{ford2023}.
\item[0.9e] X activates ERK2 in a second-order step. In the cell, \mk{5}{ERK2 is activated by cAMP}~\cite{maeda1996,laub1998,maeda2004} through the receptor cAR1; X takes the place of this cAMP signal for the reason given above (an isolated cell has no external cAMP, so a receptor-driven step could never start a pulse). ERK2$^*$ is deliberately long-lived, so it integrates X; in the calibration a first-order drive left ERK2$^*$ at 5\,\% of the pool at rest, which kept the RegA brake 87\,\% off before any spike. The second-order dependence squares the spike-to-leak contrast and can be read as a lumped ultrasensitive step; the second order in X has no direct experimental counterpart.
\item[0.9e$'$] First-order return of ERK2$^*$ to ERK2. No dephosphorylation rate or phosphatase has been measured; the rate is chosen to reproduce the measured time course of ERK2 after a cAMP stimulus, \mk{1}{a peak at 50--60\,s and a return to near-basal activity within minutes}~\cite{maeda1996,laub1998,aubry1997} (4\,min in~\cite{maeda1996}, 5--8\,min in~\cite{aubry1997}). In the Laub--Loomis circuit the only measured-type ERK2 inactivation is by PKA~\cite{laub1998} (reaction 10.D3); a PKA-independent decay is not part of it, so 0.9e$'$ is an added term. The value $\tau=50$\,s ($>\!95\,\%$ decayed after 3\,min) lies at the fast end of the reported range and was fixed by the requirement that resting ERK2$^*$ stays low.
\item[0.9f] X activates the adenylyl cyclase ACA once ERK2$^*$ has risen. ERK2 was reported to be \mk{1}{required for receptor-mediated activation of adenylyl cyclase}~\cite{segall1995}, but later work showed that \mk{3}{ERK2 is not essential for ACA activation and acts on cAMP accumulation through RegA}~\cite{maeda2004}; the ERK2$^*$ gate on ACA is therefore a design choice that orders the two events, not a measured dependence. ERK2$^*$ disarms RegA (Module~10). The ERK2$^*$ gate makes ACA follow the release of the RegA brake by a few seconds, so the cAMP made by ACA$^*$ is not destroyed immediately.
\end{description}

\paragraph{Constants.}
The parameters of 0.9a--d and 0.9y are calibrated, not measured: the excitable core has three targets (excitation threshold, spike height and recovery time), and the values in Table~\ref{tab:m0const} reproduce them.
The ignition rate $k_0$ is the only parameter that sets the firing rate.
It was calibrated on the stochastic engine, because neither a well-mixed reduction (about $2\times$ too slow: ignition nucleates in one voxel before diffusion levels it out) nor a 13-voxel deterministic reduction predicts it.
With the refractory clock active, $k_0=3.5$ molecules\,s$^{-1}$ gives a mean interval of 726\,s (95\,\% interval 706--747\,s, 8 seeds), close to the \mk{1}{12\,min} of~\cite{ford2023}; the previous value 2.52 gave 930\,s because the block adds about 390\,s of dead time.
The 12\,min of the text of~\cite{ford2023} and the mean of 5.3\,min of its fitted gamma distribution (Fig.~1E) differ; the model follows the former. \emph{Limitation:} because the block is deterministic and makes up a large part of the interval, the coefficient of variation of the isolated pacemaker is 0.16, below the value $1/\sqrt{6.6}\approx0.39$ implied by the \mk{3}{gamma distribution} measured in~\cite{ford2023}.
In populations the period converges onto the block; for 1000 cells it is $467\pm16$\,s (7.8\,min), of the order of the \mk{8}{7\,min period} of~\cite{maeda2004}.
The clock parameters $k_{Rp},k_{Rd},K_R$ were set on the measured recovery curve of the relay (no relayed pulse within 300\,s of the previous pulse, one in three at 360\,s, all by 420\,s; three seeds), not from the analytical estimate, because the block lifts when the gated relay drive crosses the excitation threshold, which depends on the relay gain.

The Hill gates $c(H)$, $b(\mathrm{Rf})$ and the ERK2$^*$ gate are justified in Table~\ref{tab:hill}.
The choice $n=4$ for the three is not a cooperativity claim; it makes the transition sharp enough that the gate is either nearly open or nearly shut at the values the input actually takes.
